\subsection*{Formal Composition Law and Selection Conditions}
\label{sec:composition-law}

We formalize the layered mappings between ontological strata $L_0\to L_1\to L_2\to L_3$ by
introducing explicit mapping objects $M_{i,j}:L_i\to L_j$ for $i<j$ and a composition operator $\circ$
defined where the codomain of the right operand equals the domain of the left operand. The central
composition requirement used throughout this manuscript is

\begin{equation}
M_{0,3} = M_{2,3}\circ M_{1,2}\circ M_{0,1}.
\end{equation}

These mappings are required to satisfy the following conditions (axioms):

\begin{enumerate}
\item \textbf{Associativity.} For indices $a<b<c<d$ we require
\[(M_{c,d}\circ M_{b,c})\circ M_{a,b} = M_{c,d}\circ (M_{b,c}\circ M_{a,b}).\]
\item \textbf{Identity maps.} For each layer $L_i$ there exists an identity mapping $\mathrm{id}_i$ such that
\[M_{i,i+1}\circ \mathrm{id}_i = M_{i,i+1},\qquad \mathrm{id}_{i+1}\circ M_{i,i+1} = M_{i,i+1}.
\]
\item \textbf{Continuity / stability under perturbation.} Let $d_i$ be a metric on the space of mappings
$L_i\to L_{i+1}$. Then small perturbations $\delta$ in $M_{i,i+1}$ (i.e. $\|\delta\|<\epsilon$) induce
controlled perturbations in composed maps. Formally, for a composition operator we assume a bound of the form
\[\|\Delta(M_{0,3})\| \le C(\epsilon),\qquad C(\epsilon)\to 0 \text{ as }\epsilon\to 0,\]
which encodes robustness to noise and finite precision.
\item \textbf{Resource-bounded realization (CPU / cost invariance).} Each mapping $M_{i,i+1}$ is realized by
resource-bounded computation with a cost functional $\mathrm{cost}_i(M_{i,i+1})$. Composition must respect
resource constraints; i.e. there exists an aggregation function $f$ such that
\[\mathrm{cost}_{0,3} \le f\big(\mathrm{cost}_{0,1},\mathrm{cost}_{1,2},\mathrm{cost}_{2,3}\big).
\]
We treat additive aggregation $f(x,y,z)=x+y+z$ or convex combinations as working hypotheses and
consider them explicit modelling choices whose empirical consequences are examined in the specs.
\item \textbf{Locality / superposition clauses.} Where mappings act linearly we assume additivity
$M(x+y)=M(x)+M(y)$. Where nonlinearity is essential we state the relevant closure properties explicitly.
\end{enumerate}

\paragraph{Remarks.}
The axioms above lift the layer-stack from a typing discipline to a compositional, category-like structure.
When stronger claims are made (for example, uniqueness of $M_{0,3}$) an additional selection principle is
required, typically in the form of an optimality condition: $M_{i,i+1}$ minimizes an explicitly defined cost
functional among feasible maps. If uniqueness cannot be proven, the space of admissible mappings should be
characterized and a selection dynamics (see Section~\ref{sec:uos-dynamics}) employed to explain observed choices.

% End composition law
